\section{Problema}

Seja a matrícula \(m=d_1d_2d_3d_4d_5d_6=318375\). Logo,
\[
(d_1,d_2,d_3,d_4,d_5,d_6)=(3,1,8,3,7,5).
\]
A substituição desses dígitos nas expressões do enunciado resulta nas matrizes
\[
A_1=\begin{bmatrix}
20&1&42&27\\
14&20&6&10\\
2&0&8&31\\
0&0&0&34
\end{bmatrix},\qquad
A_2=\begin{bmatrix}
-8&-1&0&2&1\\
-7&-2&9&7&2\\
4&3&16&-2&-2\\
-7&3&4&-9&-2\\
0&-2&-7&-7&-16
\end{bmatrix}.
\]

\subsection{Problema 1: decomposição de Schur e sensibilidade}

Para \(A_1\), deseja-se determinar \(Q\) e \(T\), com \(Q\) ortogonal e \(T\) real quase triangular superior, tais que
\[
A_1=QTQ^T.
\]
Em seguida, devem-se obter os autovalores de \(A_1\) a partir dos blocos diagonais de \(T\) e discutir a sensibilidade do problema de autovalor.

\subsection{Problema 2: autovalor dominante}

Para \(A_2\), deseja-se aproximar o autovalor dominante, isto é, o autovalor de maior módulo, pelo método da potência. O vetor inicial deve ser escolhido de modo adequado, e a iteração deve ser encerrada quando o erro relativo entre duas aproximações sucessivas do autovalor satisfizer a tolerância especificada. Deve-se também apresentar a evolução das aproximações e dos erros ao longo das iterações.
